% Input asset for encounter_multimodal_prr_v2_supplement.tex

\section{Exact prescribed finite modality: complete proof}
\label{sec:exactmfull-proof}

This section proves the exact finite-mode statement used in the main article;
Theorem~\ref{thm:exactmfull-doi} below is Theorem~1 of the main text.
The construction is pointwise in a fixed finite dimension and
a fixed finite mode count.  Its two small parameters are ordered: first the
narrow-noise parameter is fixed sufficiently small, and only then is the Doi
budget taken sufficiently small.  All stationary-point counts below are on
the declared compact positive-time window.

\subsection{Exact quotient, stationary midpoint variance, and whole-window
contact}
\label{subsec:exactmfull-physical}

Fix finite integers \(d\ge2\) and \(m\ge1\), a transverse period \(W>0\), a
contact radius \(0<a<W/2\), and a compact positive-time interval
\(I=[\tau,T]\subset(0,\infty)\).  On
\(\mathbb R\times\mathbb R\times\mathbb T_W^{d-1}\), let the longitudinal
midpoint and relative coordinates solve
\begin{align}
 \dd Z_t&=-\gamma(Z_t-\bar z)\,\dd t
       +\varepsilon\sqrt{D_0}\,\dd W_t,
 \label{eq:exactmfull-midpoint-sde}\\
 \dd R_{\parallel,t}&=-\gamma R_{\parallel,t}\,\dd t
       +2\varepsilon\sqrt{D_0}\,\dd W_{\parallel,t},
 \label{eq:exactmfull-relative-long-sde}\\
 \dd R_{\perp,t}&=2\varepsilon\sqrt{D_0}\,\dd W_{\perp,t}
       \pmod W,
 \label{eq:exactmfull-relative-perp-sde}
\end{align}
where \(D_0,\gamma>0\).  All Brownian drivers and all initial variables are
mutually independent, and
\begin{align}
 Z_0&\sim N\!\left(z_0,\varepsilon^2\frac{D_0}{2\gamma}\right),
 &z_0&\ne\bar z,
 \label{eq:exactmfull-midpoint-initial}\\
 R_{\parallel,0}&\sim
 N(r_{\parallel,0},\varepsilon^2u_0^2),
 &0<u_0^2&<\frac{4D_0}{\gamma},
 \label{eq:exactmfull-relative-long-initial}\\
 R_{\perp,0}&\sim
 \operatorname{WN}(r_{\perp,0},\varepsilon^2\Sigma_{\perp,0}),
 &\Sigma_{\perp,0}&\succ0.
 \label{eq:exactmfull-relative-perp-initial}
\end{align}
Here \(\operatorname{WN}\) denotes the wrapped Gaussian on the transverse
torus.  The two strict variance inequalities
\(D_0/(2\gamma)<D_0/\gamma\) and
\(u_0^2<4D_0/\gamma\) are precisely the unbounded Gaussian conditions that
place this initial law in the weighted state space used by
Theorem~\ref{thm:mixed-jet}; the wrapped factor is square integrable against
the uniform torus measure.

The midpoint has mean and variance
\begin{equation}
 \mu(t)=\bar z+(z_0-\bar z)\e^{-\gamma t},
 \qquad
 \operatorname{Var}Z_t=\varepsilon^2\frac{D_0}{2\gamma}
 \quad(t\ge0).
 \label{eq:exactmfull-stationary-variance}
\end{equation}
Thus only the variance, not the mean, is stationary.  This constant variance
is what produces a common-variance Gaussian mixture below.

Since \(\mu'\) has constant sign, write \(\operatorname{sgn}\mu'\in\{-1,1\}\);
fix a physical reference length \(\ell_0>0\), and define
\begin{equation}
 x(t)=\frac{\operatorname{sgn}(\mu')\,\mu(t)}{\ell_0},
 \qquad v_0:=\inf_{t\in I}x'(t)>0.
 \label{eq:exactmfull-increasing-coordinate}
\end{equation}
For the present design \(x'(t)=\gamma|z_0-\bar z|\e^{-\gamma t}/\ell_0\), so
\(v_0=\gamma|z_0-\bar z|\e^{-\gamma T}/\ell_0\); the constant \(K\) of
Lemma~\ref{lem:exactmfull-posterior-sectors} is chosen of order \(1/v_0\)
below, which is why later windows require smaller \(\varepsilon\).
Choose target times and their dimensionless centres by
\begin{equation}
 \tau<t_1<\cdots<t_m<T,
 \qquad c_j=x(t_j).
 \label{eq:exactmfull-targets}
\end{equation}
Then \(c_1<\cdots<c_m\), and the physical centre of slab \(j\) is
\(\operatorname{sgn}(\mu')\,\ell_0c_j=\mu(t_j)\).  In particular, reversing the trajectory
orientation reverses the physical centres at the same time; no centre is
held fixed under the coordinate reversal.

Write
\begin{equation}
 r_*(t)=(r_{\parallel,0}\e^{-\gamma t},r_{\perp,0})
 \label{eq:exactmfull-relative-deterministic-path}
\end{equation}
and impose the whole-window, minimum-image contact margin
\begin{equation}
 \sup_{t\in I}|r_*(t)|_{\rm mi}\le a-\eta
 \quad\text{for some }\eta>0.
 \label{eq:exactmfull-whole-window-contact}
\end{equation}
The word ``whole-window'' is essential: the condition is required on every
point of \(I\), not merely near the target times.

\begin{lemma}[Whole-window differentiated contact tail]
\label{lem:exactmfull-contact-tail}
Let
\begin{equation}
 c_{d,\varepsilon}(t)
 =\mathbb P\{|R_t|_{\rm mi}<a\}.
 \label{eq:exactmfull-contact-probability}
\end{equation}
Under Eqs.~\eqref{eq:exactmfull-relative-long-sde}--%
\eqref{eq:exactmfull-whole-window-contact}, for \(r=0,1,2\) there are finite
constants \(C_{r,d},N_{r,d},q_d>0\), independent of all sufficiently small
\(\varepsilon\), such that
\begin{equation}
 \left\|\partial_t^r(c_{d,\varepsilon}-1)\right\|_{L^\infty(I)}
 \le C_{r,d}\varepsilon^{-N_{r,d}}
       \exp(-q_d/\varepsilon^2).
 \label{eq:exactmfull-contact-tail-bound}
\end{equation}
Consequently, \(c_{d,\varepsilon}\ge1/2\) for sufficiently small
\(\varepsilon\), and the first two logarithmic derivatives of
\(c_{d,\varepsilon}\) are uniformly bounded on \(I\) and converge to zero
faster than any power of \(\varepsilon\).
\end{lemma}

\begin{proof}
In an unwrapped transverse chart the relative Gaussian has mean \(r_*(t)\)
and covariance \(\varepsilon^2\Sigma_R(t)\), with longitudinal and transverse
blocks
\begin{align}
 u_0^2\e^{-2\gamma t}
 +\frac{2D_0}{\gamma}(1-\e^{-2\gamma t}),
 \qquad
 \Sigma_{\perp,0}+4D_0t I_{d-1}.
 \label{eq:exactmfull-relative-covariance}
\end{align}
Because \(I\subset(0,\infty)\), \(u_0>0\), and
\(\Sigma_{\perp,0}\succ0\), this covariance and its first two time
derivatives are uniformly bounded on \(I\), and its eigenvalues lie between
two positive constants.  The reverse triangle inequality and
Eq.~\eqref{eq:exactmfull-whole-window-contact} give
\begin{equation}
 \inf_{t\in I}\inf_{|r|_{\rm mi}\ge a}
 d_{\rm cyl}(r,r_*(t))\ge\eta.
 \label{eq:exactmfull-contact-separation}
\end{equation}
The assumption \(a<W/2\) keeps the contact ball away from the transverse cut
locus.  Express the wrapped density as its Gaussian lattice-image sum.
Every time derivative of order at most two is the same Gaussian multiplied
by a polynomial in the scaled displacement and in \(\varepsilon^{-1}\),
with degree bounded independently of \(t\).  Equation
\eqref{eq:exactmfull-contact-separation}, the uniform covariance bounds, and
a standard Gaussian tail estimate therefore give, after summing the lattice
images,
\begin{equation}
 \sup_{t\in I}\int_{|r|_{\rm mi}\ge a}
 |\partial_t^r p_{\varepsilon,t}(r)|\,\dd r
 \le C_{r,d}\varepsilon^{-N_{r,d}}
       \e^{-q_d/\varepsilon^2},
 \qquad r=0,1,2.
 \label{eq:exactmfull-density-tail}
\end{equation}
Differentiation under the integral is justified by the same integrable
majorant and yields Eq.~\eqref{eq:exactmfull-contact-tail-bound}.  For small
\(\varepsilon\), its \(r=0\) case gives \(c_{d,\varepsilon}\ge1/2\), while
\begin{align}
 |\partial_t\log c_{d,\varepsilon}|
 &\le2|\partial_t c_{d,\varepsilon}|,\nonumber\\
 |\partial_t^2\log c_{d,\varepsilon}|
 &\le2|\partial_t^2c_{d,\varepsilon}|
      +4|\partial_tc_{d,\varepsilon}|^2.
 \label{eq:exactmfull-log-contact-bounds}
\end{align}
This proves the logarithmic-derivative statement.
\end{proof}

Fix \(\rho>0\), put
\begin{equation}
 S_*^2=\frac{D_0}{2\gamma}+\rho^2,
 \qquad \sigma=\frac{\varepsilon S_*}{\ell_0},
 \label{eq:exactmfull-common-width}
\end{equation}
and use the normalized physical slabs
\begin{equation}
 \phi_{j,\varepsilon}(z)
 =\frac{1}{\sqrt{2\pi}\,\varepsilon\rho}
  \exp\!\left[-\frac{(z-\mu(t_j))^2}
  {2\varepsilon^2\rho^2}\right].
 \label{eq:exactmfull-slabs}
\end{equation}
Let \(0<w_*\le1/m\) and
\begin{equation}
 \mathcal W_{w_*}
 =\left\{w=(w_1,\ldots,w_m):
   \sum_{j=1}^m w_j=1,\quad w_j\ge w_*>0\right\}.
 \label{eq:exactmfull-weight-set}
\end{equation}
For \(w\in\mathcal W_{w_*}\), define the Doi killing field
\begin{equation}
 K_{B,w,\varepsilon}(Z,R)
 =\mathbf{1}_{\{|R|_{\rm mi}<a\}}
  \frac{B}{W^{d-1}}\sum_{j=1}^m
  w_j\phi_{j,\varepsilon}(Z)
 =B V_{w,\varepsilon}(Z,R).
 \label{eq:exactmfull-killing}
\end{equation}
Every slab has unit full longitudinal integral, so \(B\) is the same
installed centre-space budget for every allocation.

\begin{lemma}[Exact common-variance free-exposure factorization]
\label{lem:exactmfull-factorization}
Let \(T_0(t)\) be the unkilled semigroup and \(q_0\) the initial law above.
The unit-budget free-exposure clock is exactly
\begin{align}
 G_{\varepsilon,w}(t)
 &:=\langle V_{w,\varepsilon},T_0(t)q_0\rangle
 =a_\varepsilon(t)H_{\sigma,w}(x(t)),
 \label{eq:exactmfull-factorization}\\
 a_\varepsilon(t)
 &:=\frac{c_{d,\varepsilon}(t)}
 {W^{d-1}\sqrt{2\pi}\,\varepsilon S_*},
 \label{eq:exactmfull-slow-factor}\\
 H_{\sigma,w}(x)
 &:=\sum_{j=1}^m w_j
 \exp\!\left[-\frac{(x-c_j)^2}{2\sigma^2}\right].
 \label{eq:exactmfull-mixture}
\end{align}
The positive factor \(a_\varepsilon\) has uniformly bounded first two
logarithmic time derivatives for all sufficiently small \(\varepsilon\).
\end{lemma}

\begin{proof}
Independence factors the contact event from every longitudinal slab
expectation.  By Eq.~\eqref{eq:exactmfull-stationary-variance}, convolution
of \(\phi_{j,\varepsilon}\) with the law of \(Z_t\) gives
\begin{equation}
 \mathbb E[\phi_{j,\varepsilon}(Z_t)]
 =\frac{1}{\sqrt{2\pi}\,\varepsilon S_*}
 \exp\!\left[-\frac{(\mu(t)-\mu(t_j))^2}
 {2\varepsilon^2S_*^2}\right].
 \label{eq:exactmfull-gaussian-convolution}
\end{equation}
Since
\(\mu(t)-\mu(t_j)=\operatorname{sgn}(\mu')\,\ell_0(x(t)-c_j)\),
Eqs.~\eqref{eq:exactmfull-common-width} and
\eqref{eq:exactmfull-gaussian-convolution} give
Eqs.~\eqref{eq:exactmfull-factorization}--%
\eqref{eq:exactmfull-mixture}.  The last assertion follows from
Lemma~\ref{lem:exactmfull-contact-tail}, because the remaining factors in
\(a_\varepsilon\) are independent of time.
\end{proof}

\subsection{The pure common-variance mixture}
\label{subsec:exactmfull-pure-mixture}

Set \(J=x(I)=[\alpha,\beta]\).  The targets are interior, so
\begin{equation}
 \alpha<c_1<\cdots<c_m<\beta.
 \label{eq:exactmfull-centre-order}
\end{equation}
For \(m\ge2\), let
\begin{equation}
 \Delta_j=c_{j+1}-c_j,
 \qquad \Delta_*:=\min_{1\le j<m}\Delta_j>0.
 \label{eq:exactmfull-gaps}
\end{equation}
No empty minimum is used when \(m=1\).

Define
\begin{equation}
 q_j(x)=w_j\e^{-(x-c_j)^2/(2\sigma^2)},
 \qquad
 \pi_j(x)=\frac{q_j(x)}{H_{\sigma,w}(x)},
 \qquad
 \bar c(x)=\sum_{j=1}^m\pi_j(x)c_j.
 \label{eq:exactmfull-posterior}
\end{equation}
The logarithmic slope \(L_{\sigma,w}=\partial_x\log H_{\sigma,w}\) obeys
the exact identities
\begin{align}
 L_{\sigma,w}(x)
 &=\frac{\bar c(x)-x}{\sigma^2},
 \label{eq:exactmfull-log-slope}\\
 L_{\sigma,w}'(x)
 &=\frac{\operatorname{Var}_{\pi(x)}(c)}{\sigma^4}
   -\frac{1}{\sigma^2}.
 \label{eq:exactmfull-log-slope-derivative}
\end{align}
Indeed, differentiation of \(q_j\) gives the first identity, while
\(\pi_j'=\pi_j(c_j-\bar c)/\sigma^2\) gives the second.

The mode cap for univariate Gaussian mixtures is classical
\cite{carreiraPerpinanWilliams2003modes,amendolaEngstromHaase2020modes}.
We retain the direct multiplicity-counting proof: it makes the pure-mixture
count of Theorem~\ref{thm:exactmfull-pure-topology} exhaustive including
multiple zeros, and it is the statement machine-checked in
Sec.~\ref{sec:supp-lean}.  The explicit root locations, curvature scales,
and full-complement margins that the Doi transfer needs, and that the zero
bound alone does not supply, come from
Theorem~\ref{thm:exactmfull-pure-topology} and
Lemma~\ref{lem:exactmfull-posterior-sectors} below.

\begin{lemma}[Global extended-Chebyshev zero bound]
\label{lem:exactmfull-zero-bound}
For distinct real centres and positive weights,
\(H_{\sigma,w}'\) has at most \(2m-1\) real zeros counted with
multiplicity.  Its real zero set is finite.
\end{lemma}

\begin{proof}
Multiplication by a nowhere-zero function preserves both zeros and their
multiplicities, and direct differentiation gives
\begin{equation}
 \sigma^2\e^{x^2/(2\sigma^2)}H_{\sigma,w}'(x)
 =\sum_{j=1}^m w_j(c_j-x)
   \e^{-c_j^2/(2\sigma^2)}\e^{c_jx/\sigma^2}.
 \label{eq:exactmfull-exponential-polynomial}
\end{equation}
We prove the slightly more general claim that a nonzero function
\begin{equation}
 P_m(x)=\sum_{j=1}^m(a_j+b_jx)\e^{\lambda_jx},
 \qquad \lambda_1<\cdots<\lambda_m,
 \label{eq:exactmfull-general-exp-polynomial}
\end{equation}
with no identically zero summand has at most \(2m-1\) real zeros counted
with multiplicity.

First, the zero set is finite.  After multiplication by
\(\e^{-\lambda_1x}\), the first affine term dominates as \(x\to-\infty\);
the last affine-exponential term dominates as \(x\to+\infty\).  Each
dominating affine factor has a fixed sign sufficiently far in its tail.
Thus all zeros lie in a compact interval.  A nonzero real-analytic function
has only finitely many zeros there.

Now induct on \(m\).  The \(m=1\) case is the affine zero bound.  For
\(m>1\), put \(Q=\e^{-\lambda_1x}P_m\).  Two derivatives annihilate the first
affine term.  For every \(j\ge2\), the second derivative of
\((a_j+b_jx)\e^{(\lambda_j-\lambda_1)x}\) is a nonzero affine polynomial
times the same exponential, and the remaining positive exponents are
distinct.  If a real-analytic function such as \(Q\) has \(N\ge2\) real
zeros counted with
multiplicity, generalized Rolle counting gives at least \(N-2\) zeros of
its second derivative counted with multiplicity: differentiation reduces
the multiplicity at each root, and ordinary Rolle supplies a root between
each pair of distinct roots.  Hence
\begin{equation}
 N(Q)-2\le N(Q'')\le2(m-1)-1,
 \label{eq:exactmfull-rolle-count}
\end{equation}
so \(N(P_m)\le2m-1\); the cases \(N<2\) are immediate.  In
Eq.~\eqref{eq:exactmfull-exponential-polynomial}, every affine coefficient
has nonzero slope
\(-w_j\e^{-c_j^2/(2\sigma^2)}\), so the claim applies.
\end{proof}

\begin{lemma}[Uniform adjacent isolation]
\label{lem:exactmfull-adjacent-isolation}
For the fixed finite centre set and the compact weight family, there are
\(C_{\rm iso},q_{\rm iso}>0\) such that, for every \(j<m\),
\begin{equation}
 \sup_{x\in[c_j,c_{j+1}]}
 \frac{\sum_{k\notin\{j,j+1\}}q_k(x)}
      {q_j(x)+q_{j+1}(x)}
 \le C_{\rm iso}\e^{-q_{\rm iso}/\sigma^2}
 \label{eq:exactmfull-adjacent-isolation}
\end{equation}
for all sufficiently small \(\sigma\), uniformly in \(w\).  On
\([\alpha,c_1]\), component \(1\) has the analogous one-component bound, and
on \([c_m,\beta]\), component \(m\) does.
\end{lemma}

\begin{proof}
For \(x\in[c_j,c_{j+1}]\) and \(k<j\),
\begin{equation}
 (x-c_k)^2-(x-c_j)^2
 =(c_j-c_k)(2x-c_j-c_k)\ge(c_j-c_k)^2.
 \label{eq:exactmfull-left-square-comparison}
\end{equation}
For \(k>j+1\), comparison with \(c_{j+1}\) gives the corresponding positive
lower bound.  Since \(w_k/w_j\le1/w_*\), summation over the finite centre set
proves Eq.~\eqref{eq:exactmfull-adjacent-isolation}.  The same
difference-of-squares comparison proves the two tail statements.  The
posterior mean and variance of the full mixture therefore differ from those
of the isolated adjacent pair by \(O(\e^{-q/\sigma^2})\), uniformly in the
gap, the allocation, and sufficiently small \(\sigma\).
\end{proof}

\begin{theorem}[Exact topology of the pure mixture]
\label{thm:exactmfull-pure-topology}
Uniformly for \(w\in\mathcal W_{w_*}\), every sufficiently small
\(\sigma>0\) gives exactly \(m\) simple local maxima and \(m-1\) simple local
minima of \(H_{\sigma,w}\) on \(J\), ordered alternately, with nonzero
endpoint derivatives.  The maximum near \(c_j\), denoted \(p_j(\sigma,w)\),
satisfies
\begin{equation}
 p_j=c_j+O(\e^{-q/\sigma^2}),
 \qquad L_{\sigma,w}'(p_j)
 =-\sigma^{-2}+o(\sigma^{-2}).
 \label{eq:exactmfull-peak-location}
\end{equation}
For \(m\ge2\), define the weighted equality point
\begin{equation}
 s_j(\sigma,w)
 =\frac{c_j+c_{j+1}}2
  +\frac{\sigma^2}{\Delta_j}\log\frac{w_j}{w_{j+1}}.
 \label{eq:exactmfull-weighted-crossover}
\end{equation}
The minimum in the \(j\)-th gap, denoted \(r_j(\sigma,w)\), satisfies
\begin{align}
 r_j&=s_j+O(\sigma^4),
 \label{eq:exactmfull-valley-location}\\
 L_{\sigma,w}'(r_j)
 &=\frac{\Delta_j^2}{4\sigma^4}+O(\sigma^{-2}).
 \label{eq:exactmfull-valley-curvature}
\end{align}
All remainders are uniform on the compact weight set.
\end{theorem}

\begin{proof}
Fix a constant \(A>1\).  On
\(|x-c_j|\le A\sigma^2\), one-component isolation gives
\begin{align}
 L_{\sigma,w}(x)
 &=\frac{c_j-x}{\sigma^2}
   +O(\sigma^{-2}\e^{-q/\sigma^2}),
 \label{eq:exactmfull-peak-slope-expansion}\\
 L_{\sigma,w}'(x)
 &=-\sigma^{-2}
   +O(\sigma^{-4}\e^{-q/\sigma^2}).
 \label{eq:exactmfull-peak-derivative-expansion}
\end{align}
At \(c_j-\sigma^2\) the slope is positive, at
\(c_j+\sigma^2\) it is negative, and it is strictly decreasing between
them.  Hence there is one simple maximum with the location and derivative in
Eq.~\eqref{eq:exactmfull-peak-location}.

For a gap, the isolated adjacent odds are exactly
\begin{equation}
 \frac{q_{j+1}(x)}{q_j(x)}
 =\exp\!\left[\frac{\Delta_j(x-s_j)}{\sigma^2}\right].
 \label{eq:exactmfull-adjacent-odds}
\end{equation}
Put \(h_j=(\log9)/\Delta_j\) and
\begin{equation}
 C_j=[s_j-h_j\sigma^2,s_j+h_j\sigma^2].
 \label{eq:exactmfull-crossover-interval}
\end{equation}
At the left and right endpoints the adjacent posterior masses are,
respectively, \((9/10,1/10)\) and \((1/10,9/10)\).  By
Lemma~\ref{lem:exactmfull-adjacent-isolation}, after reducing a common
\(\sigma_0\), both adjacent posterior masses are bounded below throughout
\(C_j\), and
\begin{equation}
 \operatorname{Var}_{\pi(x)}(c)\ge\kappa_0\Delta_*^2,
 \qquad
 L_{\sigma,w}'(x)\ge\frac{\kappa_v}{\sigma^4}>0
 \quad(x\in C_j)
 \label{eq:exactmfull-crossover-monotonicity}
\end{equation}
for constants independent of \(j,w,\sigma\).  At the two endpoints of
\(C_j\), Eq.~\eqref{eq:exactmfull-log-slope} gives, for sufficiently small
\(\sigma\),
\begin{equation}
 L_{\sigma,w}(s_j-h_j\sigma^2)
 \le-\frac{\Delta_j}{5\sigma^2},
 \qquad
 L_{\sigma,w}(s_j+h_j\sigma^2)
 \ge \frac{\Delta_j}{5\sigma^2}.
 \label{eq:exactmfull-crossover-signs}
\end{equation}
Thus \(C_j\) contains exactly one simple minimum root.

At \(s_j\), the adjacent posterior masses are equal and the adjacent
posterior mean is \((c_j+c_{j+1})/2\).  Hence
\begin{equation}
 L_{\sigma,w}(s_j)
 =-\frac{1}{\Delta_j}\log\frac{w_j}{w_{j+1}}
 +O(\sigma^{-2}\e^{-q/\sigma^2})=O(1)
 \label{eq:exactmfull-slope-at-crossover}
\end{equation}
uniformly in \(w\).  The mean-value theorem and
Eq.~\eqref{eq:exactmfull-crossover-monotonicity} give
\(r_j-s_j=O(\sigma^4)\).  At \(r_j\), the two adjacent posterior masses are
\(1/2+O(\sigma^2)\); substitution in
Eq.~\eqref{eq:exactmfull-log-slope-derivative} proves
Eq.~\eqref{eq:exactmfull-valley-curvature}.

For small \(\sigma\), the constructed peak and gap neighbourhoods are
pairwise disjoint and contained in \((\alpha,\beta)\).  They contain
\(m+(m-1)=2m-1\) distinct simple stationary points.  The global bound in
Lemma~\ref{lem:exactmfull-zero-bound} excludes every additional real
stationary point, including multiple ones.  Tail isolation and the fixed
separations \(c_1-\alpha>0\) and \(\beta-c_m>0\) give
\begin{equation}
 L_{\sigma,w}(\alpha)>0,
 \qquad L_{\sigma,w}(\beta)<0,
 \label{eq:exactmfull-endpoint-signs}
\end{equation}
so neither endpoint is stationary.  When \(m=1\), there are no gaps or
valley definitions: the exact identity
\(L_{\sigma,w}(x)=(c_1-x)/\sigma^2\) gives one maximum and both endpoint
signs directly.
\end{proof}

\subsection{Full posterior-sector certificate and a slow positive factor}
\label{subsec:exactmfull-slow-factor}

The zero bound above does not apply after multiplication by a nonconstant
positive factor.  We therefore certify the complete complement of small peak
and valley boxes.

\begin{lemma}[Posterior-sector certificate]
\label{lem:exactmfull-posterior-sectors}
Fix \(K>0\).  There are constants
\(A_{\rm p},A_{\rm v},C_{\rm p},C_{\rm v},
\kappa_{\rm p},\kappa_{\rm v}>0\) and \(\sigma_K>0\), depending only on the
fixed centre set, \(J\), \(w_*\), and \(K\), such that, for every
\(0<\sigma<\sigma_K\) and every admissible \(w\), the boxes
\begin{align}
 P_j&=[c_j-A_{\rm p}\sigma^2,c_j+A_{\rm p}\sigma^2],
 \label{eq:exactmfull-peak-boxes}\\
 V_j&=[r_j-A_{\rm v}\sigma^4,r_j+A_{\rm v}\sigma^4]
 \quad(j<m)
 \label{eq:exactmfull-valley-boxes}
\end{align}
are pairwise disjoint and contained in \((\alpha,\beta)\), and the following
hold:
\begin{align}
 |L_{\sigma,w}|&\le C_{\rm p},&
 L_{\sigma,w}'&\le-\kappa_{\rm p}\sigma^{-2}
 &&\text{on }P_j,
 \label{eq:exactmfull-peak-box-control}\\
 |L_{\sigma,w}|&\le C_{\rm v},&
 L_{\sigma,w}'&\ge \kappa_{\rm v}\sigma^{-4}
 &&\text{on }V_j.
 \label{eq:exactmfull-valley-box-control}
\end{align}
At the left and right box boundaries, respectively, the slope signs are
\begin{equation}
 \begin{array}{c|cc}
 &\text{left boundary}&\text{right boundary}\\ \hline
 P_j&L_{\sigma,w}\ge2K&L_{\sigma,w}\le-2K\\
 V_j&L_{\sigma,w}\le-2K&L_{\sigma,w}\ge2K.
 \end{array}
 \label{eq:exactmfull-box-boundary-signs}
\end{equation}
On the entire complement in \(J\) of the open boxes,
\begin{equation}
 |L_{\sigma,w}(x)|\ge2K,
 \label{eq:exactmfull-complement-margin}
\end{equation}
with the exhaustive alternating sign order
\begin{equation}
 +\;P_1\;-\;V_1\;+\;P_2\;-\;\cdots\;+\;P_m\;-.
 \label{eq:exactmfull-complement-sign-order}
\end{equation}
For \(m=1\), every valley clause is omitted.
\end{lemma}

\begin{proof}
Choose \(A_{\rm p}\ge4K\).  The uniform expansions
\eqref{eq:exactmfull-peak-slope-expansion}--%
\eqref{eq:exactmfull-peak-derivative-expansion} give, after reducing
\(\sigma_K\), Eq.~\eqref{eq:exactmfull-peak-box-control} and the peak row of
Eq.~\eqref{eq:exactmfull-box-boundary-signs}.  The convex-hull property
\(\bar c(x)\in[c_1,c_m]\) gives the full outer-tail estimates without any
local expansion:
\begin{align}
 x\le c_1-A_{\rm p}\sigma^2&\Longrightarrow
 L_{\sigma,w}(x)\ge A_{\rm p},\nonumber\\
 x\ge c_m+A_{\rm p}\sigma^2&\Longrightarrow
 L_{\sigma,w}(x)\le-A_{\rm p}.
 \label{eq:exactmfull-outer-tail-sector}
\end{align}

Fix a gap and abbreviate \(\Delta=\Delta_j\), \(s=s_j\).  The right
adjacent posterior mass is
\begin{equation}
 p(x)=\frac{q_{j+1}(x)}{q_j(x)+q_{j+1}(x)}
 =\frac{1}{1+\exp[-\Delta(x-s)/\sigma^2]}.
 \label{eq:exactmfull-logistic-posterior}
\end{equation}
On the crossover interval \(C_j\) in
Eq.~\eqref{eq:exactmfull-crossover-interval}, its edge odds are \(1/9\) and
\(9\), not exponentially small or large.  Both adjacent masses are therefore
bounded below throughout \(C_j\), and
Eq.~\eqref{eq:exactmfull-crossover-monotonicity} holds there.  Conversely,
the fixed finite centre diameter bounds the posterior variance, so
\(|L'|\le C\sigma^{-4}\) on \(C_j\).

The pure root satisfies \(L(r_j)=0\).  Choose \(A_{\rm v}\) large enough that
\(\kappa_vA_{\rm v}\ge2K\).  Since
\(A_{\rm v}\sigma^4=o(\sigma^2)\), the valley box lies in \(C_j\) for small
\(\sigma\).  Integrating the lower and upper bounds for \(L'\) from \(r_j\)
proves Eq.~\eqref{eq:exactmfull-valley-box-control}, the valley row of
Eq.~\eqref{eq:exactmfull-box-boundary-signs}, and the required negative and
positive bounds on \(C_j\setminus\operatorname{int}V_j\).

It remains to cover both outer sectors of every gap; this is the step that
prevents an unproved dominance claim at a crossover edge.  Consider first
\begin{equation}
 [c_j+A_{\rm p}\sigma^2,s_j-h_j\sigma^2]
 \label{eq:exactmfull-left-gap-sector}
\end{equation}
and split it at \(c_j+\Delta_j/4\), which lies in the stated interval for all
sufficiently small \(\sigma\), uniformly in \(w\).  On the portion nearer
\(c_j\), every other component is exponentially small relative to component
\(j\), so
\begin{equation}
 \bar c(x)-x\le-\tfrac12A_{\rm p}\sigma^2,
 \qquad L_{\sigma,w}(x)\le-A_{\rm p}/2\le-2K.
 \label{eq:exactmfull-left-near-peak-sector}
\end{equation}
On the remaining portion, monotonicity of
Eq.~\eqref{eq:exactmfull-logistic-posterior} and the left crossover edge give
\(p(x)\le1/10\), while \(x-c_j\ge\Delta_j/4\).  The adjacent-pair mean then
obeys
\begin{equation}
 c_j+\Delta_jp(x)-x
 \le\frac{\Delta_j}{10}-\frac{\Delta_j}{4}
 =-\frac{3\Delta_j}{20}.
 \label{eq:exactmfull-left-crossover-sector}
\end{equation}
The nonadjacent correction is exponentially small by
Lemma~\ref{lem:exactmfull-adjacent-isolation}; hence this sector also has
\(L\le-2K\) for small \(\sigma\), including its \(1/9\)-odds edge.

Symmetrically, split
\([s_j+h_j\sigma^2,c_{j+1}-A_{\rm p}\sigma^2]\) at
\(c_{j+1}-\Delta_j/4\).  On the portion nearer the crossover,
\(p(x)\ge9/10\) and \(c_{j+1}-x\ge\Delta_j/4\), giving the positive analogue
of Eq.~\eqref{eq:exactmfull-left-crossover-sector}; on the portion nearer
\(c_{j+1}\), one-component isolation gives
\(L\ge A_{\rm p}/2\ge2K\).  Thus the right outer sector has \(L\ge2K\).

The fixed centre and endpoint separations make all boxes disjoint and
interior for one common sufficiently small \(\sigma_K\).  The outer tails,
all peak boundaries, both outer sectors of every gap, and every
crossover-minus-valley sector cover every point of the box complement.
This proves Eqs.~\eqref{eq:exactmfull-complement-margin}--%
\eqref{eq:exactmfull-complement-sign-order}.  When \(m=1\), the same result
follows directly from \(L=(c_1-x)/\sigma^2\) and
Eq.~\eqref{eq:exactmfull-outer-tail-sector}.
\end{proof}

\begin{theorem}[Slow positive factors preserve the exact topology]
\label{thm:exactmfull-slow-factor}
Let \(x\in C^2(I)\) be strictly increasing with
\(\inf_Ix'=v_0>0\).  Let \(a_\sigma\in C^2(I)\) be positive and assume
\begin{equation}
 M_0:=\sup_{\sigma<\sigma_0}
 \|\partial_t\log a_\sigma\|_{L^\infty(I)}<\infty,
 \qquad
 M_1:=\sup_{\sigma<\sigma_0}
 \|\partial_t^2\log a_\sigma\|_{L^\infty(I)}<\infty.
 \label{eq:exactmfull-slow-factor-hypothesis}
\end{equation}
Then, uniformly for \(w\in\mathcal W_{w_*}\), all sufficiently small
\(\sigma>0\) give exactly \(m\) nondegenerate maxima and \(m-1\)
nondegenerate minima of
\begin{equation}
 F_{\sigma,w}(t)=a_\sigma(t)H_{\sigma,w}(x(t))
 \label{eq:exactmfull-slow-product}
\end{equation}
on \(I\), ordered alternately, with no stationary endpoint.  Relative to the
pure roots in the \(x\) coordinate,
\begin{align}
 x(t_{j,\sigma}^{\max})-p_j(\sigma,w)&=O(\sigma^2),
 \label{eq:exactmfull-slow-peak-shift}\\
 x(t_{j,\sigma}^{\min})-r_j(\sigma,w)&=O(\sigma^4),
 \label{eq:exactmfull-slow-valley-shift}
\end{align}
uniformly in \(w\).  Thus
\begin{equation}
 x(t_{j,\sigma}^{\min})=s_j(\sigma,w)+O(\sigma^4).
 \label{eq:exactmfull-slow-valley-location}
\end{equation}
\end{theorem}

\begin{proof}
Put
\begin{equation}
 b_\sigma=\partial_t\log a_\sigma,
 \qquad
 D_{\sigma,w}(t)=\partial_t\log F_{\sigma,w}(t)
 =b_\sigma(t)+x'(t)L_{\sigma,w}(x(t)).
 \label{eq:exactmfull-stationarity-function}
\end{equation}
Since \(F_{\sigma,w}>0\), its stationary points are precisely the zeros of
\(D_{\sigma,w}\), and
\begin{equation}
 D_{\sigma,w}'
 =b_\sigma'+x''L_{\sigma,w}+(x')^2L_{\sigma,w}'.
 \label{eq:exactmfull-stationarity-derivative}
\end{equation}
Choose \(K>(M_0+1)/v_0\) and apply
Lemma~\ref{lem:exactmfull-posterior-sectors}.  On the full complement of the
open boxes, the term \(x'L\) strictly dominates \(b_\sigma\), so \(D\) has
the alternating signs in Eq.~\eqref{eq:exactmfull-complement-sign-order} and
has no complement zero.

On a peak box, Eqs.~\eqref{eq:exactmfull-peak-box-control} and
\eqref{eq:exactmfull-stationarity-derivative} give
\begin{equation}
 D_{\sigma,w}'
 \le M_1+\|x''\|_\infty C_{\rm p}
 -v_0^2\kappa_{\rm p}\sigma^{-2}<0
 \label{eq:exactmfull-peak-root-monotonicity}
\end{equation}
for small \(\sigma\).  Its boundary signs give exactly one zero.  At that
zero \(F''=FD'<0\), hence the root is a nondegenerate maximum.  Similarly,
on a valley box,
\begin{equation}
 D_{\sigma,w}'
 \ge-M_1-\|x''\|_\infty C_{\rm v}
 +v_0^2\kappa_{\rm v}\sigma^{-4}>0,
 \label{eq:exactmfull-valley-root-monotonicity}
\end{equation}
so its unique zero is a nondegenerate minimum.  The certified tail signs
exclude endpoint stationary points.

At a slow-factor root, Eq.~\eqref{eq:exactmfull-stationarity-function} gives
\(|L|\le M_0/v_0\).  The pure peak root has \(L(p_j)=0\), and the uniform
negative \(L'=\Theta(\sigma^{-2})\) on its box implies the peak displacement
in Eq.~\eqref{eq:exactmfull-slow-peak-shift}.  The uniform positive
\(L'=\Theta(\sigma^{-4})\) on a valley box gives
Eq.~\eqref{eq:exactmfull-slow-valley-shift}.  Combining the latter with
Eq.~\eqref{eq:exactmfull-valley-location} proves
Eq.~\eqref{eq:exactmfull-slow-valley-location}.  Every constant is uniform
over the compact allocation family.
\end{proof}

\subsection{Fixed-\texorpdfstring{\(\varepsilon\)}{epsilon} Doi transfer and
theorem boundary}
\label{subsec:exactmfull-doi-transfer}

For \(B>0\), let
\begin{equation}
 f_{B,\varepsilon,w}(t)
 =\mathbb E_0\!\left[
 K_{B,w,\varepsilon}(Z_t,R_t)
 \exp\!\left(-\int_0^tK_{B,w,\varepsilon}(Z_s,R_s)\,\dd s\right)
 \right]
 \label{eq:exactmfull-reaction-density}
\end{equation}
be the Doi reaction-time density and set
\begin{equation}
 \mathcal F_{B,\varepsilon,w}=f_{B,\varepsilon,w}/B.
 \label{eq:exactmfull-rescaled-density}
\end{equation}
This is budget-rescaling, not normalization to unit time integral.

\begin{theorem}[Exact \(m\) Doi modes for fixed finite \((d,m)\)]
\label{thm:exactmfull-doi}
Fix all data in Eqs.~\eqref{eq:exactmfull-midpoint-sde}--%
\eqref{eq:exactmfull-killing}, including the stationary midpoint variance,
mutual independence, the strict weighted-space variance inequalities, target
times strictly inside \(I\), consistently oriented physical centres,
normalized slabs of physical width \(\varepsilon\rho\), the whole-window
contact margin, and the compact allocation family.  Then there is
\(\varepsilon_0>0\), depending on these fixed data, such that for every
\(0<\varepsilon<\varepsilon_0\) and every
\(w\in\mathcal W_{w_*}\), the continuum free-exposure clock
\(G_{\varepsilon,w}\) has exactly \(m\) nondegenerate maxima and \(m-1\)
nondegenerate minima on \(I\), ordered alternately, and no stationary
endpoint.

For each fixed such \(\varepsilon\) and allocation set
\(\mathcal A\subseteq\mathcal W_{w_*}\), define
\begin{align}
 B_{\rm top}(\varepsilon;\mathcal A)
 :=\sup\bigl\{b>0:\;&
 \mathcal F_{\beta,\varepsilon,w}
 \text{ has that complete signature on }I
 \nonumber\\[-2pt]
 &\text{for every }w\in\mathcal A
 \text{ and every }0<\beta<b\bigr\}.
 \label{eq:exactmfull-btop}
\end{align}
Then
\(B_{\rm top}^{\rm unif}(\varepsilon):=
B_{\rm top}(\varepsilon;\mathcal W_{w_*})>0\).
Consequently \(\mathcal F_{B,\varepsilon,w}\) has the stated signature for
every \(0<B<B_{\rm top}^{\rm unif}(\varepsilon)\) and every admissible
allocation.  Multiplication by \(B>0\) changes no stationary point, so the
statement also holds for the physical density \(f_{B,\varepsilon,w}\).
\end{theorem}

\begin{proof}
In Lemma~\ref{lem:exactmfull-factorization}, take
\(a_\sigma=a_\varepsilon\), with
\(\sigma=\varepsilon S_*/\ell_0\).  Lemma
\ref{lem:exactmfull-contact-tail} supplies
Eq.~\eqref{eq:exactmfull-slow-factor-hypothesis}; hence
Theorem~\ref{thm:exactmfull-slow-factor} proves the exact stationary
signature of \(G_{\varepsilon,w}\), uniformly over the compact allocation
family.

Fix one \(0<\varepsilon<\varepsilon_0\).  The ordered simple roots depend
continuously on \(w\) by the implicit-function theorem.  Their finitely many
graphs over the compact set \(\mathcal W_{w_*}\) are compact and disjoint.
Continuity of the second time derivative gives pairwise disjoint root tubes
on which \(G''\) has a uniform nonzero sign and magnitude.  After shrinking
the tubes if necessary, \(G'\) has fixed opposite signs at the two boundaries
of each tube.  On the compact complement in
\(I\times\mathcal W_{w_*}\), \(|G'|\) has a positive minimum; the two endpoint
slopes have positive uniform absolute minima as well.

For this fixed \(\varepsilon\), \(V_{w,\varepsilon}\) is bounded uniformly in
\(w\), and Eqs.~\eqref{eq:exactmfull-midpoint-initial}--%
\eqref{eq:exactmfull-relative-perp-initial} put \(q_0\) in the required
weighted state space.  Because \(w_j\ge w_*>0\), the compact allocation set
lies strictly inside the simplex and admits the affine control chart and a
complex tube required by Theorem~\ref{thm:mixed-jet}.  (If the set is the
singleton \(w_j=1/m\), the same conclusion is the corresponding pointwise
case.)  Applying that theorem with time orders \(0,1,2\) and no control
derivative gives
\begin{equation}
 \sup_{w\in\mathcal W_{w_*}}
 \|\mathcal F_{B,\varepsilon,w}-G_{\varepsilon,w}\|_{C^2(I)}
 \longrightarrow0
 \qquad(B\downarrow0).
 \label{eq:exactmfull-fixed-epsilon-c2-transfer}
\end{equation}
Choose \(b_\varepsilon>0\) so that the \(C^1\) error is smaller than all
boundary, complement, and endpoint margins and the second-derivative error
is smaller than the root-tube curvature margin.  In each tube,
\(\mathcal F'\) has opposite boundary signs and \(\mathcal F''\) keeps the
sign of \(G''\), so it has exactly one root of the same type.  The complement
and endpoints contain none.  Thus \(b_\varepsilon\) belongs to the defining
set in Eq.~\eqref{eq:exactmfull-btop}, so
\(B_{\rm top}^{\rm unif}(\varepsilon)\ge b_\varepsilon>0\).  This proves the
complete finite-window signature.
\end{proof}

\begin{remark}[Quantifiers, saturation, and scope]
\label{rem:exactmfull-scope}
The quantifier order is
\begin{equation}
 \text{fix }(d,m,\text{all geometric and probabilistic data})
 \ \Longrightarrow\ 
 0<\varepsilon<\varepsilon_0
 \ \Longrightarrow\ 
 0<B<B_{\rm top}^{\rm unif}(\varepsilon).
 \label{eq:exactmfull-quantifier-order}
\end{equation}
There is no asserted lower bound for
\(B_{\rm top}^{\rm unif}(\varepsilon)\) uniform as
\(\varepsilon\downarrow0\), and no interchange of the two limits.  The
construction is not uniform in \(d\) or \(m\), does not provide one geometry
with arbitrarily many modes, does not count stationary points outside \(I\),
and gives no process-level event-mass floor.  It treats normalized
longitudinal Gaussian slabs, not arbitrary localized catalyst patches.

Moreover, Lemma~\ref{lem:exactmfull-contact-tail} shows that the relative
contact factor is asymptotically saturated on the theorem window:
\begin{equation}
 c_{d,\varepsilon}=1
 +O(\varepsilon^{-N}\e^{-q/\varepsilon^2})
 \quad\text{in }C^2(I).
 \label{eq:exactmfull-saturated-contact}
\end{equation}
Thus the exact Doi embedding is rigorous, but the modal basis is created by
the ordered narrow slabs sampled by the monotone constant-variance midpoint;
the relative encounter factor is a slow perturbation in this asymptotic
family.  Nontrivial contact, a useful common positive budget, deterministic
finite-parameter complement certificates, survival, and event-basin masses
remain separate continuum-validation requirements rather than consequences
of Theorem~\ref{thm:exactmfull-doi}.
\end{remark}
